\subsection{積分の比}

公式は積分の値を積で表したものである.今度は逆に,積分の大小を使ってその積の極限を求める.まず$I_n$の値を並べてみよう.
\begin{center}
\begin{tabular}{cccccccccc}
\toprule
$n$ & $0$ & $1$ & $2$ & $3$ & $4$ & $5$ & $6$ & $7$ & $8$\\
\midrule
$I_n$ & $1.5708$ & $1$ & $0.7854$ & $0.6667$ & $0.5890$ & $0.5333$ & $0.4909$ & $0.4571$ & $0.4295$\\
\bottomrule
\end{tabular}
\end{center}
$I_n$は減少し,しかも隣り合う値の差が次第に小さくなっていく.隣どうしの比$\frac{I_{n+1}}{I_n}$は$1$へ近づくだろう.

実際,\,$0\le\sin x\le1$だから,指数を大きくすると被積分関数は小さくなる.したがって
\[
 I_{n+2}\le I_{n+1}\le I_n
\]
である.各辺を$I_n>0$で割り,式\eqref{式:ウォリス漸化式}を使うと
\[
 \frac{n+1}{n+2}
 =\frac{I_{n+2}}{I_n}
 \le\frac{I_{n+1}}{I_n}\le1
\]
となる.両端が$1$へ近づくため
\begin{equation}\label{式:ウォリス隣接比}
 \lim_{n\to\infty}\frac{I_{n+1}}{I_n}=1
\end{equation}
である.各$I_n$の極限を先に求めなくても,比の極限は分かるのだ.

ここで,次の有限の積を作る.
\[
 P_m=\prod_{k=1}^{m}\frac{(2k)^2}{(2k-1)(2k+1)}
\]
$\prod$は,指定された範囲の項をすべて掛ける記号である.最初の数項を書けば
\[
 P_m=\frac{2\cdot2}{1\cdot3}\cdot
      \frac{4\cdot4}{3\cdot5}\cdots
      \frac{2m\cdot2m}{(2m-1)(2m+1)}
\]
となる.分子の偶数の積は$\{(2m)!!\}^2$,分母の奇数の積は$(2m-1)!!\,(2m+1)!!$である.一方,式\eqref{式:ウォリス偶数}と式\eqref{式:ウォリス奇数}の比を取ると
\begin{align*}
 \frac{I_{2m+1}}{I_{2m}}
 &=\frac{2}{\pi}\cdot
   \frac{\{(2m)!!\}^2}{(2m-1)!!\,(2m+1)!!}\\
 &=\frac{2}{\pi}P_m
\end{align*}
となる.左辺は$1$へ近づくので,次の極限が得られる.
\begin{equation}\label{式:ウォリスの積}
 \boxed{
 \lim_{m\to\infty}\prod_{k=1}^{m}
 \frac{4k^2}{4k^2-1}=\frac{\pi}{2}}
\end{equation}
これが\textbf{ウォリスの積}である.無限個の数を一度に掛けたと考えるのではなく,有限の積$P_m$の数列の極限として定義している.

\paragraphbf{円周率を挟む}
証明の途中の不等式を残しておけば,近似の誤差も評価できる.\,$n=2m$とすると
\[
 \frac{2m+1}{2m+2}
 \le\frac{I_{2m+1}}{I_{2m}}\le1
\]
だから
\[
 \frac{2m+1}{2m+2}\cdot\frac{\pi}{2}
 \le P_m\le\frac{\pi}{2}
\]
である.各辺に正の数を掛けて$\pi$について解くと
\[
 2P_m\le\pi\le\frac{2m+2}{2m+1}\cdot2P_m
\]
となる.実際に計算してみよう.
\begin{center}
\begin{tabular}{ccc}
\toprule
$m$ & $2P_m$ & $\frac{2m+2}{2m+1}\cdot2P_m$\\
\midrule
$1$ & $2.6667$ & $3.5556$\\
$10$ & $3.0677$ & $3.2138$\\
$100$ & $3.1338$ & $3.1494$\\
$1000$ & $3.1408$ & $3.1424$\\
\bottomrule
\end{tabular}
\end{center}
$\pi=3.14159\ldots$は,どの$m$でも左右の値の間にある.有限の計算の結果が,どれほど確かな範囲を与えるかまで分かったのだ.ただし,区間の幅は$\frac{2P_m}{2m+1}$であり,\,$\frac{\pi}{2m+1}$より小さいものの,\,$m$に反比例する程度でしか縮まらない.\,$m=1000$でも,確定するのは$3.14$までである.
